\chapter{Statistics for digital data}

\section{Random variables}
%{{{

A random variable is a mathematical formalization of a quantity or
object which depends on random events. Specifically, a random variable
$X$ is a deterministic function that maps an \emph{outcome}
$\omega\in\Omega$ to a \emph{realization} $x\in\mathcal{X}$:
\begin{equation}
  X:\Omega\rightarrow \mathcal{X}.
  \label{eq:DRV}
\end{equation}
$\Omega$ is known as the \emph{sample space} and $\mathcal{X}$ is
called the \emph{realization space}. Notice the randomness comes from
the fact that the underlying outcome is random. $\mathcal{X}$ must be
also measurable (otherwise it could be impossible to assign a
probability, such as for example, $P(X\leq x)$).

\begin{itemize}
\item An example of a random variable is the The 10-Minutes Bus
  Scenario. Imagine you arrive at a bus stop where a bus comes exactly
  every 10 minutes. You don't know when the last one left, so you
  could wait anywhere from 0 to 10 minutes. In this example, the
  outcome $\omega$ is the actual, physical event of you standing at
  the curb and the bus eventually pulling up, and the sample space
  $\Omega$ is continuous because you could arrive at any point in time
  between the bus cycles. The random variable $X$ is ``how many
  minutes you wait'', and $x\in\mathbb{R}$ is such number. Obviously
  $X$ is measurable, because we can compute, for example,
  $P(X\leq 5)$. This is an example of a continuous random variable.
\end{itemize}

%}}}
  
\section{\glsentryfullpl{DRV}} % Discrete Random Variable
%{{{

A \gls{DRV} is a random variable where $\mathcal{X}$ is conuntable,
regardless of whether the input $\Omega$ is discrete or
continuous.\footnote{It is easy to recognize a \gls{DRV} becayse they
  usually involve ``enumerating'' something. For example, imagine you
  want to report if you hit the bullseye in a dart board (with ``1''
  if the dart land on the bullseye, and ``0'' otherwise). The input
  (sample space, i.e., the exact location where the dart hits) is a
  continuous and uncountable (an infinite decimal coordinate)
  space. Notice that even though at the input we are dealing with an
  infinite, uncountable number of points where the dart could land,
  the random variable only has two possible outputs:
  $\{0,1\}$. Continuous random variables usually involve ``measuring''
  something. Examples: The exact height of a person, the time it takes
  for a battery to die, or the temperature outside.}

\begin{itemize}
\item Example (continuous input): Let's consider again the 10-Minutes
  Bus Scenario. Here $\Omega=[0,10[$. Now, if we define
  $\mathcal{X}=\{0,1\}$ as:
\begin{enumerate}
\item $X(\omega)=1$ if you wait more than 5 minutes ($\omega<5$), and
\item $X(\omega)=0$ otherwise,
\end{enumerate}
$X$ would be (continuous to) discrete.

\item Example (discrete input): If we toss a coin, and define
  $\mathcal{X}=\{0,1\}$ as:
\begin{enumerate}
\item $X(\omega)=0$ if $\omega=\text{``Heads''}$, and
\item $X(\omega)=1$ if $\omega=\text{``Tails''}$,
\end{enumerate}
$X$ would be (discrete to) discrete.
\end{itemize}

%}}}

\section{\glsentryfullpl{DRVec}} % Discrete random vector
%{{{

A \gls{DRVec} (also called a \emph{multivariate random variable}) is a
generatization of a \gls{DRV} where $\mathcal{X}$ is
$N$-dimensional. For this reason, when we deal with \glspl{DRVec}, we
will denote the sample space with $\bm{\Omega}$, a sample as
$\bm{\omega}$, the \gls{DRVec} as $X$, a realization with $\bm{x}$,
and the realization space as $\mathcal{X}^N$.

\begin{itemize}
\item For example, suppose we roll two ($N=2$) fair six-sided
  dice. The sample space $\Omega$ consists of 36 ordered pairs
  $(d_1,d_2)$ where $d_i\in\{1,\cdots,6\}$. Now, we define a
  \gls{DRVec} $X=[S,D]^{\text{T}}$, where $S=d_1 + d_2$ and
  $D=|d_1-d_2|$ are \glspl{DRV}. For example, if
  $\bm{\omega}=[4,6]^{\text{T}}$, we get
  $\bm{x}=X(\bm{\omega})=[10,2]^{\text{T}}$.
  
  Notice that the vector representation expresses a dependency between
  $S$ and $D$. For example, if you know hat $S=12$, then the $D$ must
  be 0 (the dice must be 6 and 6). Such dependency allows the
  definition of the joint \gls{PMF}, and the existence of the concepts
  of covariance and correlation.
\end{itemize}

%}}}

\section{A statistical model of signal digitalization}
%{{{

\label{sec:statistical_model}
The digitalization of a sequence of $N$ signal-samples can be modeled
statistically using a \gls{DRVec}.\footnote{By physical requirements,
  digital signals are finite in length.} We can identify that:
\begin{enumerate}
\item $\bm{\Omega}$ (the sample space) is the set of all possible physical
  continuous-time waveforms that the sensor could encounter during the
  sampling window, and $\bm{\omega}$ (the ``analog reality'', i.e., the
  signal to digitalize) is a single, specific continuous voltage
  waveform that actually occurs during your measurement
  window.
\item $X$ represents the \gls{ADC}. It is a deterministic
  function that performs two operations on the outcome $\bm{\omega}$:
  \begin{enumerate}
  \item \emph{Sampling}: $X$ looks at $\bm{\omega}$ at $N$
    specific time points $\{t_1,t_2,\cdots,t_N\}$.
  \item \emph{Quantization}: Then $X$ usually rounds the voltage at those
    points to the nearest integer.
  \end{enumerate}
\item The realization $\bm{x}=X(\bm{\omega})$ is the digizalized
  signal. Notice that $X$ is a vector operator that inputs a
  sample(-segment) (a window of signal) and outputs a sequence of
  digital numbers $\{\bm{x}_1,\bm{x}_2,\cdots,\bm{x}_N\}$.
\item $\mathcal{X}^{N}$ (the realization space) is the set of all possible
  realizations (all posible digital signals of length $N$).
\end{enumerate}

%}}}

\section{\acrlongpl{SRV}}
%{{{

A \gls{SRV} is the result of a stochastic process whose statistical
properties do not change. For example, if we say that a digital signal
is a \gls{SRV}, then the probability distribution of a signal at any
given set of time/space points is the same as the distribution at
those same samples shifted by any constant amount in time or space.

%}}}

\section{Expectation (mean) of a \gls{DRV}}
%{{{

The mean, expected value, expectation or expectancy of a random
variable $X$, denoted by $\mathbb{E}[X]$ or simply by $\mu$, is its
long-term \emph{average} outcome. It is defined as the center of mass
of the probability distribution. Specifically,
\begin{equation}
  \mathbb{E}[X] = \sum_{x\in\mathcal{X}}x\Pr(X=x).
  \label{eq:expectation}
\end{equation}

The expectation is a linear operator and satisfies that
\begin{equation}
  \mathbb{E}[aX+bY] = a\mathbb{E}[X] + b\mathbb{E}[Y],
\end{equation}
even if $X$ and $Y$ are dependent.

%}}}

\section{$n$-th moment of a \gls{DRV}}
%{{{

The $n$-th moment of a \gls{DRV} $X$ is the expectation of the
$n$-th power of $X$:
\begin{equation}
  \mathbb{E}[X^n] = \sum_\Omega(X(\omega))^n\Pr(\omega).
  \label{eq:moments}
\end{equation}

%}}}

\section{The sample mean (arithmetic average)}
%{{{

\label{sec:average}
The \emph{sample mean} or \emph{arithmetic average} of the sequence of
realizations $\{x^{(i)}\}_{i=1}^{N}=\bm{x}=\{\bm{x}_i\}_{i=1}^{N}$
generated after $N$ observations of a random variable is defined by
\begin{equation}
  \overline{\bm{x}} = \frac{1}{N}\sum_{i=1}^{N}\bm{x}_i.
  \label{eq:average}
\end{equation}

Notice that by the \gls{LLN} (see Section~\ref{sec:LLN}), as the
number of realizations $N\rightarrow\infty$, the sample mean
$\overline{x}$ converges to the mean $\mu$. Therefore, the sample mean
is an estimation of the mean (the expected value) that should improve
with $N$.

%Similarly, the average of a vector(-realization)
%$\mathbf{x}$ can be found with
%\begin{equation}
%  \overline{\mathbf{x}} = \frac{1}{N}\sum_{i=1}^{N}\mathbf{x}_i,
%  \label{eq:vector_average}
%\end{equation}
%where $\mathbf{x}_i$ is the $i$-th component(-realization) of
%$\mathbf{x}$.

%}}}

\section{The \acrlong{LLN}} % Law Large Numbers
%{{{
\label{sec:LLN}
In general terms, the \gls{LLN} expresses the fact that our empirical
observations will eventually align with our theoretical models if we
collect enough data. In other words, (see Eq.~\ref{eq:average}) it is
satisfied that
\begin{equation}
  \lim_{N\rightarrow\infty}\overline{\bm{x}}=\mathbb{E}[X],
  \label{eq:LLN}
\end{equation}
where $\bm{x}=\{\bm{x}_i\}_{i=1}^{N}$ is a sequence of realizations
generated by the \gls{DRV} $X$.

%Notice that Eq.~\ref{eq:LLM} also holds for a sequence of \gls{idd}
%random variables $\{X^{(i)}\}_{i=1}^{N}$ if they satisfy that
%\begin{equation*}
%  \mathbb{E}[X^{(i)}]=\mu\quad\forall i.
%\end{equation*}

%}}}

\section{Variance}
%{{{

\label{sec:variance}
The variance $\text{Var}(X)$ tells you how far the actual outcomes
generated by a random variable $X$ are likely to deviate from that
expected value. It is defined as the expected value of the squared
deviation from the mean,
\begin{equation}
  \text{Var}(X) = \mathbb{E}[(X - \mathbb{E}[X])^2] = \mathbb{E}[X^2]-\mathbb{E}[X]^2.
  \label{eq:variance}
\end{equation}

%}}}

\section{The sample variance}
%{{{

The sample variance can be computed over the realizations $\bm{x}=\{x^{(i)}\}_{i=1}^{N}$ using
\begin{equation}
  \text{Var}(\bm{x}) = \frac{1}{N-1}\sum_{i=0}^{N}(x^{(i)}-\overline{\bm{x}})^2.
\end{equation}

%}}}

%%%%%%%%%%%%%

\section{\glsentryfull{PMF}} % Probability mass function
%{{{

The \gls{PMF}, shows the probability for each specific value that a
\gls{DRV} can take.\footnote{This information in the case of a
  continuous random variable is represented by a Probability Density
  Function (PDF).} The \gls{PMF} of a random variable (or vector) $X$
is defined as
\begin{equation}
  p(x) = \Pr(X=x), \quad x\in\mathcal{X},
\end{equation}
that is the probability with which a random variable $X$ takes
the value $x$ in a countable set $\mathcal{X}$, where it must be true that
\begin{equation*}
  p(x)\ge 0,\quad\sum_{x\in\mathcal{X}}p(x)=1.
\end{equation*}

%}}}

\section{Marginal of a random variable}
%{{{

The marginal of a variable is its probability distribution, ignoring
the other variables. For example, imagine we survey a neighborhood and
record the type of pet (cat or dog) and its color (brown or
black). The marginal of pet type considers just the number of cat and
dogs, not their color.

%}}}

\section{Conditional distribution}
%{{{

The conditional probability mass function is
$$ p(y \mid x) = \frac{p(x,y)}{p(x)} \quad \text{if } p(x) > 0, $$
where $p(x) = \sum_{y} p(x,y)$ is the marginal of $X$.

The conditional expectation of $h(x,Y)$ given $X=x$ is then
$$ E_{Y \mid X=x}[h(x,Y)] = \sum_{y} h(x,y) p(y \mid x). $$

%}}}

\section{Joint \gls{PMF}} % Probability Mass Function
%{{{

A random vector variable allows us to define a single function
$P(\mathbf{X}=\mathbf{x})$, which is shorthand for:
\begin{equation}
P(S = s \text{ and } D = d)
\end{equation}
This joint \gls{PMF} tells us the probability of specific coordinates
in the 2D space X. Without the vector framework, we could only
describe the marginal probabilities (the probability of the sum
regardless of the difference), losing the context of the experiment.

%}}}

\section{Power of a signal}
%{{{

Notice that using the expectation operator, the power of a digital
signal (see Eq.~\ref{eq:power_discrete_signal}) can be found as
\begin{equation}
  P(\mathbf{x}) = \mathbb{E}(|\mathbf{x}|^2).
  \label{eq:power_as_expectation}
\end{equation}

%}}}

\section{Expectation of the function of a discrete random variable}
%{{{

\begin{equation}
  \mathbb{E}[f(\mathbf{x})]=\sum_if(\mathbf{x}_i)\Pr(\mathbf{x}=\mathbf{x}_i).
  \label{eq:expectation_of_function}
\end{equation}

\section*{Example:}
Suppose we want to compute the expected winnings of a simple betting game tossing a coin, where if Heads, you win \$2, but if Tails, you lose \$1. We define the payoff as a function of $X$:

\begin{equation*}
  f(X)=2X-1
\end{equation*}

Therefore, if $X=1$ (Heads): $f(1)=2(1)-1=1$ (dollar net gain), but if $X=0$ (Tails): $f(0)=2(0)-1=-1$ (dollar loss). The expectation of the payoff is

\begin{equation*}
\mathbb{E}[f(X)]=f(1)\Pr(X=1)+f(0)\Pr(X=0) = (1)(12)+(-1)(12)=0.
\end{equation*}

So your expected profit in this game is zero.

%}}}

\section{Conditional distribution}
%{{{

The conditional probability mass function is
$$ p(y \mid x) = \frac{p(x,y)}{p(x)} \quad \text{if } p(x) > 0, $$
where $p(x) = \sum_{y} p(x,y)$ is the marginal of $X$.

The conditional expectation of $h(x,Y)$ given $X=x$ is then
$$ E_{Y \mid X=x}[h(x,Y)] = \sum_{y} h(x,y) p(y \mid x). $$

%}}}

\section{Law of total expectation}
%{{{

The expectation with respect to the joint distribution is
\begin{equation}
  \mathbb{E}_{(X,Y)}[h(X,Y)]=\mathbb{E}_{X}[\mathbb{E}_{Y\mid X}[h(X,Y)]].
\end{equation}


\subsection*{Example: Coin + Die}

Toss a fair coin and roll a fair die simultaneously.

\subsection*{Random Variables}
\[
C = 
\begin{cases}
1 & \text{if Heads} \\
0 & \text{if Tails}
\end{cases}, 
\qquad
D \in \{1,2,3,4,5,6\}.
\]

The coin and die are independent, so the joint distribution is
\[
p(c,d) = p(c)\,p(d) = \frac{1}{2} \cdot \frac{1}{6} = \frac{1}{12}.
\]

\subsection*{Function of Interest}
Suppose the payoff is
\[
h(C,D) = C \cdot D,
\]
i.e., you get the die roll if the coin shows Heads, and 0 if Tails.

\subsection*{Expectation w.r.t. the joint distribution}
\[
\mathbb{E}_{(C,D)}[h(C,D)] 
= \sum_{c=0}^{1} \sum_{d=1}^{6} h(c,d) \, p(c,d).
\]

%}}}

\section{Power of a signal as a function of its mean and variance}
%{{{

The average power of a random variable is equal to its variance plus its mean
(expectation) squared
\begin{equation}
  P(\mathbf{x}) =  \mathbb{V}(\mathbf{x}) + \mathbb{E}(\mathbf{x})^2.
  \label{eq:power_mean_variance}
\end{equation}
This is because
\begin{align*}
  \mathbb{E}(\mathbf{x}^2)
  & = \mathbb{E}(\mathbf{x}^2 - \mathbb{E}(\mathbf{x}) + \mathbb{E}(\mathbf{x})) \\
  & = \mathbb{E}((\mathbf{x} - \mathbb{E}(\mathbf{x}))^2 + 2\mathbb{E}(\mathbf{x})(\mathbf{x}-\mathbb{E}(\mathbf{x})) + \mathbb{E}(\mathbf{x})^2) \\
  & = \mathbb{E}((\mathbf{x}-\mathbb{E}(\mathbf{x}))^2) + 2\mathbb{E}(\mathbf{x})(\mathbb{E}(\mathbf{x})-\mathbb{E}(\mathbf{x})) + \mathbb{E}(\mathbf{x})^2) \\
  & = \mathbb{V}(\mathbf{x}) + 0 + \mathbb{E}(\mathbf{x})^2.
\end{align*}

%}}}

\section{Covariance}
%{{{
\label{sec:covariance}

The covariance $\text{Cov}(\bm{x}, \bm{y})$ (also denoted by
$R_{\bm{x}\bm{y}}$) is a measure of how the realizations of two random
variables $\bm{x}$ and $\bm{y}$ change together. In simpler terms, it
tells us the direction of the linear relationship between two
variables. The covariance between two discrete signals $\bm{x}$
and $\bm{y}$ is calculated as
%\begin{equation}
%  \text{Cov}(\textbf{x}, \textbf{y}) = \mathbb{E}[(\mathbf{x}-\overline{\mathbf{x}})(\mathbf{y}-\overline{\mathbf{y}})].
%\end{equation}
%\begin{equation}
%  \mathbb{V}(\mathbf{x}) = \mathbb{E}\left((\mathbf{x} - \mathbb{E}(\mathbf{x}))^2 \right).
%\end{equation}
\begin{equation}
  \text{Cov}(\bm{x}, \bm{y}) = \mathbb{E}\big((\bm{x}-\mathbb{E}(\bm{x}))(\bm{y}-\bm{E}(\bm{y}))\big).
\end{equation}

Notice that (see Eq.~\ref{eq:variance})
\begin{equation}
  \text{Cov}(\bm{x}, \bm{x}) = \mathbb{E}[\bm{x}\bm{x}^{\text{T}}]= \text{Var}(\bm{x})\bm{I},
\end{equation}
where $\bm{I}$ is the identity matrix.

%}}}

\section{Covariance matrix}
%{{{
\label{sec:covariance_matrix}

The covariance matrix $\Sigma_{\overrightarrow{\mathbf{x}}}$ of a
random vector
$\overrightarrow{\mathbf{x}}=[\mathbf{x}_1,\cdots,\mathbf{x}_N]^T$,
defined as
\begin{equation}
  (\Sigma_{\overrightarrow{\mathbf{x}}})_{i,j}=\mathbb{C}(\mathbf{x}_i,\mathbf{x}_j),
\end{equation}
is a $N\times N$ matrix
\begin{equation}
\Sigma_{\overrightarrow{\mathbf{x}}} = 
\begin{pmatrix}
\mathbb{V}(\mathbf{x}_1) & \mathbb{C}(\mathbf{x}_1, \mathbf{x}_2) & \cdots & \mathbb{C}(\mathbf{x}_1, \mathbf{x}_p) \\
\mathbb{C}(\mathbf{x}_2, \mathbf{x}_1) & \mathbb{V}(\mathbf{x}_2) & \cdots & \mathbb{C}(\mathbf{x}_2, \mathbf{x}_p) \\
\vdots & \vdots & \ddots & \vdots \\
\mathbb{C}(\mathbf{x}_p, \mathbf{x}_1) & \mathbb{C}(\mathbf{x}_p, \mathbf{x}_2) & \cdots & \mathbb{V}(\mathbf{x}_p)
\end{pmatrix}
\end{equation}
that express the covariance between the random variables of a random vector.

%}}}

\section{$L_2$ norm}
\label{sec:L2_norm}
%{{{
\label{sec:covariance_matrix}

$L_2$ norm (also called \emph{magnitude} and \emph{Euclidean norm}) of
a discrete signal $\mathbf{x}$ is defined by
\begin{equation}
  ||\mathbf{x}||_2 = \sqrt{\sum_i\mathbf{x}_i^2}.
\end{equation}
Notice that the L2 norm and the Mean Square Error (MSE) are closely
related concepts, because
\begin{equation}
  ||\mathbf{x} - \mathbf{y}||_2^2 = N\cdot\text{MSE}(\mathbf{x} - \mathbf{y}),
\end{equation}
where $N$ is the length of $\mathbf{x}$.

%}}}

\section{Uniform distribution}
%{{{

The uniform distribution, usually denoted by $\mathcal{U}(a,b)$, is a
continous probability distribution in which all outcomes are equally
likely, and can be described by \gls{PDF}
\begin{equation}
  \Pr(X{=}x)) =
  \begin{cases}
    \frac{1}{b-a} & \text{for}\quad a \le x \le b, \\
    0 & \text{otherwise},
  \end{cases}
\end{equation}
where $X$ represents a continuous random variable, and $[a, b]$ is the
range of amplitudes of the random samples.

As can be seen, the mean
\begin{equation}
  \mu = \frac{a+b}{2},
\end{equation}
and the variance
\begin{equation}
  \sigma^2 = \frac{(b-a)^{2}}{12}.
\end{equation}

In Python, the method
\href{https://numpy.org/doc/stable/reference/random/generated/numpy.random.uniform.html}{\text{numpy.random.uniform}$(a=0,
  b=1, N=\text{None})$} draws sequences of realizations $\bm{x}$ from
an uniform distribution $\mathcal{U}(a,b)$. Formally, we can say that
each $\bm{x}_i\in\mathcal{X}$, being $\mathcal{X}$ the realization
space of the \gls{DRV} $X\sim\mathcal{U}(a,b)$.

%}}}

\section{Gaussian distribution}
%{{{
\label{sec:gaussian_distribution}

A continuous random variable $X$ that is generated following a Gaussian
distribution (also known as the normal distribution and usually denoted
by $\mathcal{N}(\mu, \sigma^2)$) is described by a \gls{PDF}
\begin{equation}
  \Pr(X{=}x) = \frac{1}{\sqrt{2\pi}\sigma} e^{-\frac{(x-\mu)^2}{2\sigma^2} },
  \label{eq:normal_PDF}
\end{equation}
where $x$ is a random sample of $X$, $\Pr(X{=}x)$ represents the
probability of finding $x$ in $X$, $\mu$ is the mean of $X$, and
$\sigma$ is the standard deviation of $X$.

When we work with a digital signal (quantized discrete random
variable) $\mathbf{x}$, the \gls{PDF} becomes a
\gls{PMF}\footnote{Also called discrete \gls{PDF}.}
\begin{equation}
  \Pr(\mathbf{x}{=}\mathbf{x}_i) = \frac{1}{\sqrt{2\pi}\sigma} e^{-\frac{(\mathbf{x}_i-\mu)^2}{2\sigma^2} },
  \label{eq:normal_PMD}
\end{equation}
where $\mathbf{x}_i$ is the $i$-th sample of $\mathbf{x}$.

%}}}

\section{Poisson distribution}
\label{sec:poisson_distribution}
%{{{

If a \gls{DRV} $X$ follows a Poisson distribution,
we denote it as
\begin{equation}
  X\sim\mathcal{P}(\lambda),
\end{equation}
where $\lambda$ is the \emph{shape parameter}, representing the mean
of the realizations ($\mathbb{E}[X]=\lambda$) in a specified
interval. $\lambda$ also is the variance ($\text{Var}(X)=\lambda$).

The Poisson distribution is a discrete probability distribution that
models the number of times an event occurs in a fixed
interval in the events domain usually time or space),
under the following assumptions:
\begin{enumerate}
\item \textbf{Independency}: The occurrence of one event does not
  affect the probability of another (e.g., one photon hitting a
  detector doesn't "push" another away).
\item \textbf{Constant rate}: The average rate $\lambda$ (events per
  interval) is constant.
\item \textbf{Singularity}: Two events cannot occur at the exact same
  instant.
\end{enumerate}

Its \gls{PMF} is:
\begin{equation}
  \Pr(X{=}k) = \frac{e^{-\lambda}\lambda^{k}}{k!},
  \label{eq:PD}
\end{equation}
where $\lambda$ is the mean and the variance. In the context of
microscopy, it describes the probability of a specific number of
photons or electrons hitting a pixel during an exposure.

\subsection{Normal approximation to the Poisson distribution}
As the mean $\lambda$ of a Poisson distribution increases, the
distribution becomes more symmetric and begins to resemble the bell
curve of a Gaussian (Normal) distribution. If you have a Poisson
variable $X\sim\text{Poisson}(\lambda)$, then
\begin{equation}
  X\approx\mathcal{N}(\lambda,\lambda),
  \label{eq:poisson_to_normal}
\end{equation}
and this approximation improves with $\lambda$ (in general,
$\text{Possion}(\lambda)$ and $\mathcal{N}(\lambda,\lambda)$ are
considered indistinguishable).
%}}}

\section{The \glsentryfull{CLT}}
%{{{
\label{sec:CLT}

The \acrshort{CLT} states that when a set of independent \glspl{SRV}
(see Section~\ref{sec:DSRV}) are added together, the \gls{SRV}
resulting of their normalized sum tends toward a normal (Gaussian)
distribution, even if the original variables themselves are not
normally distributed. Let
$\mathbf{n}=\{\mathbf{n}^{(i)}\}_{i=0}^{N-1}$ such set of independent
and identically distributed (i.i.d.) \glspl{SRV}, then the
\acrshort{CLT} asserts that the mean signal
\begin{equation}
  \overline{n} = \frac{1}{N}\sum_{i=0}^{N-1}\mathbf{n}^{(i)}\sim\mathcal{N}(\mu,\frac{\sigma}{N}),
\end{equation}
where
\begin{equation}
  \mathbb{E}(\mathbf{n}^{(i)})=\mu\quad\forall i,
\end{equation}
and
\begin{equation}
  \mathbb{V}(\mathbf{n}^{(i)})=\sigma\quad\forall i.
\end{equation}
Notice that
\begin{equation}
  \lim_{N\rightarrow\infty}\mathbb{V}(\overline{\mathbf{n}}) = 0.
\end{equation}

%}}}

